\documentclass[10pt,a4paper]{amsart}
\usepackage[margin=19mm]{geometry}
\usepackage{amsmath,amssymb,mathtools}
\usepackage[scr=rsfs,cal=euler]{mathalfa}
\usepackage{xfrac}
\usepackage[unicode,hidelinks,bookmarks=false]{hyperref}
\newcommand{\ie}{i.e.\ }
\newcommand{\cf}{cf.\ }
\newcommand{\resp}{respectively\ }
\newcommand{\Subsection}{\subsection}
\providecommand{\nobreakdash}{\nobreak\hbox{-}}
\setlength{\emergencystretch}{2em}
\multlinegap0pt
\usepackage{amssymb}
\usepackage{enumerate}
\usepackage{tikz}
\newtheorem{proposition}{Proposition}
\newcommand{\End}{\operatorname{End}}
\title{Conic modules, secondary fans and non-commutative resolutions}
\author{Aimeric Malter}
\date{Source: arXiv:2603.23945v1 (CC BY 4.0)}
\begin{document}
\maketitle

\noindent\emph{Source and licence.} Conic modules, secondary fans and non-commutative resolutions. Version arXiv:2603.23945v1 (CC BY 4.0), distributed by arXiv under the Creative Commons Attribution 4.0 International licence, http://creativecommons.org/licenses/by/4.0/, as recorded in the arXiv licence metadata for this exact version and shown on the abstract page https://arxiv.org/abs/2603.23945v1. Full licence notice: \texttt{licenses/arxiv-2603.23945-CC-BY-4.0.txt}.\par\smallskip

\noindent\textit{Editorial scope.} Counted unit: the article's proposition Prop:SpecialCasesWithoutNCCR (source0.tex lines 1735--1744). The article's complete original proof of it is delivered with it (source0.tex lines 1745--1791, one \texttt{proof} environment of 47 lines). No mathematics is written, repaired or re-derived: the statement and the proof are the article's own lines, in the article's own notation, and no retained line is substituted or re-flowed.  No further source-local context is needed: every cross-reference of every retained line resolves inside the two delivered spans.  Nothing was omitted from any retained span, so the package carries no omission record.  No retained line contains a citation, so the wrapper carries no bibliography and no bibliography entry.  No retained line contains a figure environment, an image-inclusion command or drawing code, so no image is delivered and the package declares no asset dependency.  Editorial formatting only, in addition: an AMS \texttt{amsart} preamble that restates the article's own declarations and macros used by the retained lines, namely the environments \texttt{proposition} on separate counters, together with the standard packages those lines need.  The retained environments print in this document as Proposition 1 in order of appearance, whereas the article prints the same items with the numbering of its own full document.  The complete native source of the article as distributed in the arXiv e-print for this exact version is bundled unchanged as \texttt{sources/197121/source0.tex}.
\begin{proposition}
    \label{Prop:SpecialCasesWithoutNCCR}
    Let $\sigma$ be an almost simplicial Gorenstein cone such that either of the following is true:
    \begin{enumerate}
        \item The collection of $\beta_\rho$ associated to $\sigma$ is, up to flipping all signs, $\{2,1,1,-2,-2\}$.
        \item The collection of $\beta_\rho$ associated to $\sigma$ is, up to flipping all signs, $\{2,2,2,-3,-3\}$.
        \item The collection of $\beta_\rho$ associated to $\sigma$ is, up to flipping all signs, $\{2,2,-1,-1,-1,-1\}$.
    \end{enumerate}
    Then there does not exist an incredulous set of conic modules, and thus there is no incomplete direct sum $\mathbb{B}$ of conic modules such that $\End_{R_\sigma}(\mathbb{B})$ is an NCCR of $R_\sigma$.
\end{proposition}

\begin{proof}
    \underline{(1):} The zonotope is $(-4,4)$, with the seven associated complexes being:
    \begin{align*}
        A_1\longrightarrow A_{-1}^{\oplus 2}\longrightarrow & A_{-3},\\
        A_1^{\oplus 2}\longrightarrow A_{-1}^{\oplus 4}\oplus A_2\longrightarrow A_{-3}^{\oplus 2}\oplus A_0^{\oplus 2}\longrightarrow & A_{-2},\\
        A_1 \longrightarrow A_{-1}^{\oplus 2}\oplus A_1\oplus A_2^{\oplus 2}\longrightarrow A_{-3}\oplus A_{-1}^{\oplus 2}\oplus A_0^{\oplus 4}\oplus A_3\longrightarrow A_{-3}\oplus A_{-2}^{\oplus 2}\oplus A_1^{\oplus 2}\longrightarrow & A_{-1},\\
        A_0\longrightarrow A_{-1}^{\oplus 4}\oplus A_0^{\oplus 2}\longrightarrow A_{-3}^{\oplus 2}\oplus A_{-2}\oplus A_0^{\oplus 2}\oplus A_1^{\oplus 4}\longrightarrow A_{-2}\oplus A_{-1}^{\oplus 2}\oplus A_{2}^{\oplus 2}\longrightarrow & A_0,\\
        A_{-1}^{\oplus 2}\longrightarrow A_{-3}\oplus A_0^{\oplus 4}\oplus A_1^{\oplus 2}\longrightarrow A_{-2}^{\oplus 2}\oplus A_{-1}\oplus A_1^{\oplus 2}\oplus A_2^{\oplus 4}\longrightarrow A_{-1}\oplus A_0^{\oplus 2}\oplus A_3^{\oplus 2}\longrightarrow & A_1,\\
        A_{-2}\longrightarrow A_{-1}^{\oplus 2}\oplus A_0\longrightarrow A_0\oplus A_1^{\oplus 2}\longrightarrow & A_2,\\
        A_{-1}\longrightarrow A_0^{\oplus 2}\oplus A_1\longrightarrow A_1\oplus A_2^{\oplus 2}\longrightarrow &  A_3.
    \end{align*}


    Observe that no subset is incredulous, hence no NCCR is obtained as endomorphism algebra of an incomplete direct sum of conic modules. 

    \underline{(2):} In this case, the zonotope is $(-6,6)$, with the 11 complexes associated to the lattice points being:

    \begin{align*}
        A_{1}\longrightarrow A_{-2}^{\oplus 2}\longrightarrow & A_{-5},\\
        A_2\longrightarrow A_{-1}^{\oplus 2}\longrightarrow & A_{-4},\\
        A_1^{\oplus 3}\longrightarrow A_{-2}^{\oplus 6}\oplus A_3\longrightarrow A_0^{\oplus 2}\oplus A_{-5}^{\oplus 3}\longrightarrow & A_{-3},\\
        A_2^{\oplus 3}\longrightarrow A_{-1}^{\oplus 6}\oplus A_4\longrightarrow A_1^{\oplus 2}\oplus A_{-4}^{\oplus 3}\longrightarrow & A_{-2},\\
        A_1^{\oplus 3}\longrightarrow A_{-2}^{\oplus 6}\oplus A_3^{\oplus 3}\longrightarrow A_{-5}^{\oplus 3}\oplus A_0^{\oplus 6}\oplus A_5\longrightarrow A_2^{\oplus 2}\oplus A_{-3}^{\oplus 3}\longrightarrow & A_{-1},\\
        A_0\longrightarrow A_{-1}^{\oplus 6}\longrightarrow A_{-4}^{\oplus 3}\oplus A_1^{\oplus 6}\longrightarrow A_3^{\oplus 2}\oplus A_{-2}^{\oplus 3}\longrightarrow & A_0,\\
        A_{-2}^{\oplus 2}\longrightarrow A_{-5}\oplus A_0^{\oplus 6}\longrightarrow A_{-3}^{\oplus 3}\oplus A_2^{\oplus 6}\longrightarrow A_4^{\oplus 2}\oplus A_{-1}^{\oplus3}\longrightarrow & A_1,\\
        A_{-1}^{\oplus 2}\longrightarrow A_{-4}\oplus A_1^{\oplus 6}\longrightarrow A_{-2}^{\oplus 3}\oplus A_3^{\oplus 6}\longrightarrow A_5^{\oplus 2}\oplus A_0^{\oplus 3}\longrightarrow & A_2,\\
        A_{-3}\longrightarrow A_{-1}^{\oplus 3}\longrightarrow A_1^{\oplus 3}\longrightarrow & A_3,\\
        A_{-2}\longrightarrow A_0^{\oplus 3}\longrightarrow A_2^{\oplus 3}\longrightarrow & A_4,\\
        A_{-1}\longrightarrow A_1^{\oplus 3}\longrightarrow A_3^{\oplus 3}\longrightarrow & A_5.
    \end{align*}


    No incredulous subset exists and hence no NCCR can be obtained via conic modules.
    
    \underline{(3):} The zonotope here is $(-4,4)$ and the associated complexes are:
    \begin{align*}
        A_1\longrightarrow A_0^{\oplus 4}\longrightarrow A_{-1}^{\oplus 6}\longrightarrow A_{-2}^{\oplus 4}\longrightarrow &A_{-3},\\
        A_2\longrightarrow A_1^{\oplus 4}\longrightarrow A_0^{\oplus 6}\longrightarrow A_{-1}^{\oplus 4}\longrightarrow& A_{-2},\\
        A_1^{\oplus 2}\longrightarrow A_0^{\oplus 8}\oplus A_3\longrightarrow A_{-1}^{\oplus 12}\oplus A_2^{\oplus 4}\longrightarrow A_{-2}^{\oplus 8}\oplus A_1^{\oplus 6}\longrightarrow A_{-3}^{\oplus 2}\oplus A_0^{\oplus 4}\longrightarrow & A_{-1},\\
        A_0\longrightarrow A_1^{\oplus 8}\longrightarrow A_0^{\oplus 12}\oplus A_3^{\oplus 4}\longrightarrow A_{-1}^{\oplus 8}\oplus A_2^{\oplus 6} \longrightarrow A_{-2}^{\oplus 2}\oplus A_1^{\oplus 4}\longrightarrow & A_0,\\
        A_{-1}^{\oplus 6}\longrightarrow A_{-2}^{\oplus 4}\oplus A_1^{\oplus 12}\longrightarrow A_{-3}\oplus A_0^{\oplus 8}\oplus A_3^{\oplus 6}\longrightarrow A_{-1}^{\oplus 2}\oplus A_2^{\oplus 4}\longrightarrow & A_1,\\
        A_{-1}^{\oplus 4}\longrightarrow A_{-2} \oplus A_1^{\oplus 8}\longrightarrow A_0^{\oplus 2}\oplus A_3^{\oplus 4}\longrightarrow & A_2,\\
        A_{-1}\longrightarrow A_1^{\oplus 2}\longrightarrow&  A_3.
    \end{align*}

    Again, no incredulous subset exists, hence we have no NCCR via conic modules.
\end{proof}

\end{document}
